% !TeX program = xelatex
\documentclass[12pt,a4paper]{article}
\usepackage[margin=2cm]{geometry}
\usepackage{fontspec}
\usepackage{amsmath,amssymb,enumitem}
\setmainfont{Times New Roman}
\setlength{\parindent}{0pt}
\setlength{\parskip}{6pt}
\begin{document}
\begin{center}
{\Large\bfseries Đáp án Toán vào lớp 10 Hà Nội năm 2020 -- phần Đại số}
\end{center}

\textbf{Bài I.}
\begin{enumerate}[label=\arabic*)]
  \item Thay $x=4$ vào $A$, ta được
  \[
  A=\frac{\sqrt{4}+1}{\sqrt{4}+2}=\frac{3}{4}.
  \]
  \item Với $x\ge 0$, $x\ne1$, ta có
  \begin{align*}
  B&=\frac{3}{\sqrt{x}-1}-\frac{\sqrt{x}+5}{x-1}\\
   &=\frac{3(\sqrt{x}+1)-(\sqrt{x}+5)}{(\sqrt{x}-1)(\sqrt{x}+1)}\\
   &=\frac{2(\sqrt{x}-1)}{(\sqrt{x}-1)(\sqrt{x}+1)}
    =\frac{2}{\sqrt{x}+1}.
  \end{align*}
  \item Từ câu 2,
  \[
  P=2\cdot\frac{\sqrt{x}+1}{\sqrt{x}+2}\cdot\frac{2}{\sqrt{x}+1}+\sqrt{x}
   =\frac{4}{\sqrt{x}+2}+\sqrt{x}.
  \]
  Do đó
  \[
  P-2=\frac{4}{\sqrt{x}+2}+\sqrt{x}-2
      =\frac{x}{\sqrt{x}+2}\ge 0.
  \]
  Dấu bằng xảy ra khi và chỉ khi $x=0$. Vậy giá trị nhỏ nhất của $P$ là $2$, đạt được tại $x=0$.
\end{enumerate}

\textbf{Bài II.}
\begin{enumerate}[label=\arabic*)]
  \item Gọi vận tốc đi bộ của An là $x$ km/h ($x>0$). Vận tốc đi xe đạp là $x+9$ km/h. Thời gian đi bộ và đi xe đạp lần lượt là $\dfrac{3}{x}$ giờ và $\dfrac{3}{x+9}$ giờ. Vì thời gian đi xe đạp ít hơn $45$ phút, tức $\dfrac{3}{4}$ giờ, nên
  \[
  \frac{3}{x}-\frac{3}{x+9}=\frac{3}{4}.
  \]
  Rút gọn phương trình:
  \[
  x^2+9x-36=0\Longleftrightarrow (x-3)(x+12)=0.
  \]
  Vì $x>0$, suy ra $x=3$. Vậy vận tốc đi bộ của An là $3$ km/h.
  \item Diện tích mặt cầu bán kính $R=2$ cm là
  \[
  S=4\pi R^2=4\cdot3{,}14\cdot2^2=50{,}24\text{ cm}^2.
  \]
\end{enumerate}

\textbf{Bài III.}
\begin{enumerate}[label=\arabic*)]
  \item Đặt $t=\dfrac{1}{y-1}$ (điều kiện $y\ne1$). Hệ trở thành
  \[
  \begin{cases}
  2x+3t=5,\\
  4x-t=3.
  \end{cases}
  \]
  Từ phương trình thứ hai, $t=4x-3$. Thay vào phương trình thứ nhất, ta được $14x=14$, nên $x=1$ và $t=1$. Do $\dfrac{1}{y-1}=1$, suy ra $y=2$. Hệ có nghiệm duy nhất $(x;y)=(1;2)$.
  \item
  \begin{enumerate}[label=\alph*)]
    \item Trên trục $Oy$ có $x=0$. Thay vào $y=mx+4$, được $y=4$. Vậy $A(0;4)$.
    \item Giao điểm của $(d)$ với trục $Ox$ có tung độ bằng $0$, nên $B\left(-\dfrac{4}{m};0\right)$. Tam giác $OAB$ vuông tại $O$. Để tam giác cân, hai cạnh góc vuông phải bằng nhau:
    \[
    OA=OB\Longleftrightarrow 4=\left|\frac{4}{m}\right|
    \Longleftrightarrow |m|=1.
    \]
    Vì $m\ne0$, các giá trị cần tìm là $m=1$ hoặc $m=-1$.
  \end{enumerate}
\end{enumerate}

\textbf{Bài V.} Điều kiện xác định là $x\ge\dfrac{2}{3}$. Ta biến đổi phương trình:
\begin{align*}
\sqrt{x}+\sqrt{3x-2}=x^2+1
&\Longleftrightarrow 2\sqrt{x}+2\sqrt{3x-2}=2x^2+2\\
&\Longleftrightarrow 2(x-1)^2+(\sqrt{x}-1)^2+(\sqrt{3x-2}-1)^2=0.
\end{align*}
Vế trái là tổng các số không âm, nên tổng bằng $0$ khi và chỉ khi từng số hạng bằng $0$. Suy ra $x=1$, giá trị này thỏa điều kiện xác định. Vậy phương trình có nghiệm duy nhất $x=1$.
\end{document}
